% Source: https://artofproblemsolving.com/community/c3h1676965_expected_valuecmms_problem_3570
% Stable identity: aops-topic-1676965-q01
Let $V$ be the expected value of the final coin. The first coin returned is a nickel, dime or quarter with probability $\dfrac{1}{4}$ each, and the process stops; with probability $\dfrac{1}{4}$ it is a penny, and then we are back where we started, so the expected final value is again $V$. Hence
\[
V=\dfrac{1}{4}\cdot5+\dfrac{1}{4}\cdot10+\dfrac{1}{4}\cdot25+\dfrac{1}{4}V\quad\Longrightarrow\quad\dfrac{3}{4}V=10\quad\Longrightarrow\quad V=\boxed{\dfrac{40}{3}\text{ cents}}\approx13.3\text{ cents}.
\]
(The process stops with probability $1$, since the chance of $n$ pennies in a row is $\left(\dfrac{1}{4}\right)^n\to0$. Equivalently, the final coin is a nickel, dime or quarter with probability $\dfrac{1}{3}$ each, giving $\dfrac{5+10+25}{3}=\dfrac{40}{3}$.) Since you put in one penny of your own, the expected profit is $\dfrac{40}{3}-1=\dfrac{37}{3}$ cents per penny you start with.
